At the full volume budget, increasing the TWoS count from two to twelve raises mean cooling from $37.58\%$ to $54.33\%$, while the radius decreases from $0.0721$ to $0.0397$. At twelve cavities, the voxel and random means are $55.02\%$ and $38.78\%$, respectively. The most useful allocation comparison is across budgets: twelve small TWoS cavities at half volume yield $47.26\%$ cooling (design range $47.03$--$47.67\%$), versus $47.45\%$ ($46.78$--$48.22\%$) for six larger cavities at full volume. The former also has $20.6\%$ less cavity surface area, by exact sphere geometry. This is similar measured cooling with less removal, not a proved equivalence or an advantage unique to TWoS. The voxel selector shows a similar allocation trend. A fourfold thinner WoS termination shell, $1.25\times10^{-4}$, changes cooling by at most $0.204$ percentage points for the pre-specified first-seed full-budget endpoint designs from both TD selectors. This sensitivity check does not bound the remaining bias.
An additional first-seed check of the half-volume comparison, chosen after inspecting the allocation results, gives $46.88\%$ versus $47.21\%$ cooling with the thinner shell, respectively.
